\documentclass[conference]{IEEEtran}

\usepackage{cite}
\usepackage{amsmath,amssymb,amsfonts}
\usepackage{algorithmic}
\usepackage{graphicx}
\usepackage{textcomp}
\usepackage{xcolor}

\def\BibTeX{{\rm B\kern-.05em{\sc i\kern-.025em b}\kern-.08em
    T\kern-.1667em\lower.7ex\hbox{E}\kern-.125emX}}

\begin{document}

\title{Analysis of Credit Fraud\\
{\footnotesize \textsuperscript{}}}

\author{\IEEEauthorblockN{Madison Byrd}
\IEEEauthorblockA{\textit{Artificial Intelligence Institute} \\
\textit{University of South Carolina}\\
Columbia, South Carolina \\
mjadebyrd@gmail.com}}

\maketitle

\begin{abstract}
Credit fraud datasets are extremely imbalanced, which can lead to issues when designing models to 
predict fraud. Creating a model to predict credit fraud is a balance between a low threshold that 
flags too many transactions and a high threshold that flags too few transactions.
\end{abstract}

\begin{IEEEkeywords}
data analysis, credit fraud, predictive models
\end{IEEEkeywords}

\section{Introduction}

\section{Related Work}
\cite{b1}

\section{Methodology}

\section{Results}

\section{Discussion}

\section{Conclusion}

\begin{thebibliography}{00}
\bibitem{b1} Author, "Title," \textit{Journal}, year.
\end{thebibliography}

\end{document}
